% Generated from locale/tr/content/first-order-logic/axiomatic-deduction/proof-theoretic-notions.tex; do not hand-edit.


\iftag{FOL}
      {\olfileid[tr]{fol}{axd}{ptn}}
      {\olfileid[tr]{pl}{axd}{ptn}}

\olsection{İspat Kuramsal Kavramlar}

\begin{explain}
Bir dizi önemli semantik kavramı
(\iftag{FOL}{geçerlilik}{totoloji}, mantıksal gerektirme, sağlanabilirlik)
tanımladığımız gibi, şimdi bunlara karşılık gelen \emph{ispat kuramsal
kavramları} tanımlıyoruz. Bunlar, !!{sentence}s ile
\iftag{FOL}{!!{structure}s}{değerlemeler} arasındaki sağlama bağıntısına değil,
belirli formüllerin !!{derivability} ya da
!!{nonderivability} durumuna dayanarak tanımlanır. Bu kavramların çakıştığının bulunması
önemliydi. Çakıştıklarını söyleyen sonuçlar \emph{sağlamlık} ve
\emph{tamlık teoremleridir}.
\end{explain}

\begin{defn}[!!^{derivability}]
Bir !!^a{formula}~$!A$, $\Gamma$'dan \emph{!!{derivable}} sayılır; bunun
koşulu $!A$ ile biten ve $\Gamma$'dan başlayan !!a{derivation} bulunmasıdır.
Bunu $\Gamma \Proves !A$ ile gösteririz.
\end{defn}

\begin{defn}[Teoremler]
Bir !!^a{formula}~$!A$, boş kümeden $!A$ için bir !!a{derivation} varsa
\emph{teorem}dir. $!A$ teoremse $\Proves !A$, değilse $\Proves/ !A$ yazarız.
\end{defn}

\begin{defn}[Tutarlılık]
!!{formula}s içeren bir $\Gamma$ kümesi ancak ve ancak
$\Gamma\Proves/ \lfalse$ ise \emph{tutarlı}, aksi hâlde \emph{tutarsız}dır.
\end{defn}

\begin{prop}[Yansımalılık]
\ollabel{prop:reflexivity}
$!A \in \Gamma$ ise $\Gamma \Proves !A$.
\end{prop}

\begin{proof}
  Tek başına !!{formula}~$!A$, $\Gamma$'dan $!A$ için bir
  !!a{derivation} oluşturur.
\end{proof}

\begin{prop}[Monotonluk]
\ollabel{prop:monotonicity}
$\Gamma \subseteq \Delta$ ve $\Gamma \Proves !A$ ise $\Delta \Proves !A$.
\end{prop}

\begin{proof}
$\Gamma$'dan $!A$ için her !!{derivation}, $\Delta$'dan $!A$ için de bir
!!a{derivation} oluşturur.
\end{proof}

\begin{prop}[Geçişlilik]
\ollabel{prop:transitivity}
$\Gamma \Proves !A$ ve $\{!A\} \cup \Delta \Proves !B$ ise
$\Gamma \cup \Delta \Proves !B$.
\end{prop}

\begin{proof}
  $\{!A\} \cup \Delta \Proves !B$ olduğunu varsayalım. O zaman
  $\{!A\} \cup \Delta$'dan $!B$ için $!B_1$, \dots, $!B_l = !B$
  biçiminde bir !!a{derivation} vardır. Bu !!{derivation} içindeki bazı
  adımlar, önceki bir $!B_i = !A$ satırına başvuran bir kuralla doğru
  olabilir. Varsayım gereği $\Gamma$'dan $!A$ için bir !!a{derivation}
  vardır; başka deyişle şu !!a{derivation} söz konusudur:
  $!A_1$, \dots, $!A_k = !A$. Burada her $!A_i$ aksiyom, $\Gamma$ içinde
  yer alan bir !!a{element} veya bir çıkarım kuralınca doğru bir adımdır.
  Şu diziyi düşünelim:
  \[
  !A_1, \dots, !A_k = !A, !B_1, \dots, !B_l = !B.
  \]
  İkinci !!{derivation} içinde $\{!A\}$ varsayımıyla gerekçelendirilen her
  $!B_i=!A$ satırını silip bu satıra yapılan başvuruları daha önceki
  $!A_k=!A$ satırına yönlendirirsek, $\Gamma\cup\Delta$'dan $!B$ için doğru
  bir türetim elde ederiz. Diğer bütün satırların gerekçeleri aynen
  kalır.
\end{proof}

Özellikle $\Gamma \Proves !A$ ve $!A \Proves !B$ ise $\Gamma \Proves !B$
olur. Ayrıca $!A_1, \dots, !A_n \Proves !B$ ve her $i$ için
$\Gamma \Proves !A_i$ ise $\Gamma \Proves !B$ sonucu çıkar.

\begin{prop}
\ollabel{prop:incons}
$\Gamma$ ancak ve ancak her $!A$ için $\Gamma \Proves !A$ olduğunda
tutarsızdır.
\end{prop}

\begin{proof}
Alıştırma.
\end{proof}

\tagprob{FOL}
\begin{prob}
\olref[fol][axd][ptn]{prop:incons} önermesini ispatlayın.
\end{prob}
\tagendprob

\tagprob{notFOL}
\begin{prob}
\olref[pl][axd][ptn]{prop:incons} önermesini ispatlayın.
\end{prob}
\tagendprob

\begin{prop}[Kompaktlık]
\ollabel{prop:proves-compact}
  \begin{enumerate}
  \item $\Gamma \Proves !A$ ise $\Gamma_0 \Proves !A$ olacak biçimde sonlu
    bir $\Gamma_0 \subseteq \Gamma$ vardır.
  \item $\Gamma$'nın her sonlu alt kümesi tutarlıysa $\Gamma$ tutarlıdır.
  \end{enumerate}
\end{prop}

\begin{proof}
  \begin{enumerate}
    \item $\Gamma \Proves !A$ ise sonlu bir !!{formula}s dizisi
      $!A_1$, \dots,~$!A_n$ vardır; $!A \ident !A_n$ olur ve her $!A_i$
      mantıksal aksiyom, $\Gamma$ içinde yer alan bir !!a{element} ya da
      önceki !!{formula}s üzerine uygulanan izinli bir çıkarım kuralının
      sonucudur.
      $\Gamma_0$, bu $!A_i$'ler arasından $\Gamma$'da bulunanların kümesi
      olsun. Aynı !!{derivation}, $\Gamma_0$ bakımından da geçerlidir; yani
      yine !!a{derivation} olur;
      dolayısıyla $\Gamma_0 \Proves !A$.
    \item Bu, $!A \ident \lfalse$ özel durumunda (1)'in karşıt tersidir.
  \end{enumerate}
\end{proof}

